\section{\env: A Scalable Tool-Integrated Agentic Training Environment for VLMs}
\label{sec:env}

\begin{figure*}[t]
    \centering
    % \vspace{-1ex}
    \includegraphics[width=\linewidth]{figure/overview.pdf}
    \caption{Overview of \env{}. \env{} contains a comprehensive suite of reasoning-intensive VQA tasks and tools in an interactive execution environment, scaling visual-centric tool-integrated agentic training for VLM agents. }
    % \vspace{-1ex}
    \label{fig:overview}
\end{figure*}

\subsection{Diverse Task and Tool Collection} \label{sec:4.1}

\noindent\textbf{Reasoning-Intensive VQA Tasks.}
\env{} comprises a unified training and evaluation environment that couples diverse multimodal tasks with executable tools and verifiers, emphasizing not only final answers but also \emph{auditable sequences of tool calls} that ground those answers. We curate \emph{verifiable} instances from 13 established benchmarks to ensure broad coverage of difficulty and reasoning types, spanning perception (vision), symbolic manipulation (math/geometry), and language understanding (document/chart interpretation). Specifically, training tasks span seven complementary axes (Appendix~\ref{app:data}): 
(1) \emph{Chart Understanding} (FigureQA~\cite{ebrahimi2018figureqa}, ChartQA~\cite{masry2022chartqa}),
(2) \emph{Geometric Reasoning} (Geometry3K~\cite{lu2021inter}, GeoQA~\cite{chen2021geoqa}, UniGeo~\cite{chen2022unigeo}),
(3) \emph{Geospatial Reasoning} (MapQA~\cite{chang2022mapqa}, InfographicVQA~\cite{mathew2022infographicvqa}),
(4) \emph{Scientific Reasoning} (ScienceQA~\cite{lu2022learn}, VizWiz~\cite{bigham2010vizwiz}),
(5) \emph{Document Understanding} (DocVQA~\cite{mathew2021docvqa}),
(6) \emph{Spatial/Compositional Reasoning} (CLEVR~\cite{johnson2017clevr}),
and (7) \emph{Others} (ThinkVL~\cite{chitty2025imagenet}, A-OKVQA~\cite{schwenk2022okvqa}).


\noindent \textbf{Visual-Centric Tool Sets.}
Small open-source VLMs often lack the fine-grained perception and domain routines (\eg, precise localization, symbolic formalization, dense chart parsing) required by the above tasks. 
\env{} therefore exposes a standardized, extensible tool interface that lets agents offload these subproblems to reliable modules, \emph{grounding} high-level reasoning in structured intermediate results (\eg, tables, text). 
Twenty-six tools are organized into four families (Appendix~\ref{app:tool}): 
(1) \emph{Perception}, such as GroundingDINO~\cite{liu2023grounding}, SAM~\cite{kirillov2023segment}, EasyOCR;
(2) \emph{Chart Understanding}, such as ChartMoE~\cite{xu2025chartmoe};
(3) \emph{Diagram Formalization}, such as CDL, Inter‑GPS;
and (4) \emph{Math Solvers}, such as G‑LLaVA~\cite{gao2025gllava}, MultiMath~\cite{peng2024multimath}.

\subsection{Executable Interactive Environment} \label{sec:4.2}

\noindent \textbf{Interface.}
\env{} exposes a Gymnasium-style API with \texttt{reset()} and \texttt{step()}.
Given an instruction $x\in\mathcal{I}$, \texttt{reset} returns the initial partial observation $o_0$ (question text and supporting image[s]) and initializes the interaction history $c_0$.
Each episode is a POMDP mentioned in \cref{sec:pre}. 
% at time step $t$, the environment yields the current observation $o_t$.
% Conditioned on the history $c_t=(x,g_1,a_1,o_1,\ldots,g_t,a_t,o_t)$, the agent then generates a reasoning token sequence $g_{t+1}$ followed by a tool call $a_{t+1}$, i.e., $(g_{t+1},a_{t+1})$ at step $t{+}1$.
% The environment executes $a_{t+1}$ and appends the resulting feedback to form $c_{t+1}$.



% \env employs a standardized environment interface that adheres closely to the widely used OpenAI Gymnasium format, featuring primary functions \texttt{.reset()} and \texttt{.step()}.
% From the agent's perspective, the environment centers on a sampled task $x\in\mathcal{I}$, initially presenting only the textual question and supportive image.
% Each interaction between the environment and the agent can be formalized as a POMDP process as detailed in Section~\ref{sec:pre}.
% At the time step $t$, the environment yields the current observation $o_t$.
% Based on the interaction history $c_t$, the agent then generates the subsequent reasoning thought and tool-call action $(g_{t+1},a_{t+1})$ at step $t+1$, contingent upon the specific task and environment design.

\noindent \textbf{Action Space.}
The action space $\mathcal{A}$ is strictly constrained by the available toolset.
We formalize an action as a typed tuple $a_t \in \mathcal{A}$ consisting of the unique identifier and the corresponding arguments passed to the interface for execution.

\noindent \textbf{Observation Space.}
The observation space $\mathcal{O}$ comprises environmental feedback after tool execution.
Each $o_t\in\mathcal{O}$ encapsulates either successful execution results or runtime error messages (if any), serving as external signals for verifying intermediate hypotheses and adjusting the subsequent reasoning trajectory.

% The observation space $\mathcal{O}$ comprises the direct feedback received from the environment following a tool invocation. 
% This observation $o_t$ encapsulates both successful execution results and runtime error messages, serving as the critical external signal for the agent to verify and adjust its reasoning trajectory.

\subsection{Scalable Training Facility} \label{sec:4.3}
\noindent \textbf{VLMs as Tools.}
To enhance visual perception, we include compute‑intensive VLMs like {G-LLaVA}~\cite{gao2025gllava} and {ChartMoE}~\cite{xu2025chartmoe} as interactive services.
To integrate them efficiently into distributed RL, we deploy a highly concurrent microservice architecture that encapsulates each VLM as an independent HTTP service with three layers: 
(i) a \textbf{FastAPI} front end that exposes RESTful endpoints with asynchronous batched requests;
(ii) an intermediate \textbf{Tool} layer that parses instruction actions, retrieves image data from trajectory metadata, and formats observations; and
(iii) a \textbf{Ray Actor} layer that keeps model weights resident in GPU memory after initialization, eliminating the prohibitive latency of repeated model (re)loading under high‑frequency tool calls.


% Among the aforementioned selected tools, G-LLaVA~\cite{gao2025gllava} and ChartMoE~\cite{xu2025chartmoe} are heavy-weight VLMs as interactive tools, and require significant amount of computational resource.
% To address the issue and integrate them into our distributed RL framework, we developed a highly concurrent microservices architecture, which wraps each VLM tool as an independent HTTP service: 
% (1) a \textbf{FastAPI} top layer that provides RESTful interfaces for asynchronous request batching;
% (2) an intermediate \textbf{Tool} layer responsible for parsing instruction actions, extracting image data from trajectory metadata, and formatting observations;
% and (3) a foundational \textbf{Ray Actor} layer, which maintains model parameters persistently in GPU memory upon initialization, eliminating the prohibitive latency overhead of repetitive model reloading during high-frequency tool calls.

\noindent \textbf{Asynchronous Training.}
To sustain high throughput during RL rollouts, we use Ray to orchestrate concurrency. At each step, the policy emits a \texttt{<think>} segment followed by a \texttt{<tool\_call>}; on generation of the sentinel token \texttt{</tool\_call>}, decoding halts and the framework assembles batched HTTP requests that include trajectory identifiers and image paths. Ray manages request queues and load balancing across tool servers. For resource efficiency, compute‑heavy VLM tools are pinned to dedicated GPUs, while lightweight utilities are multiplexed on shared CPUs.  

% To ensure high throughput for efficient large-scale RL training, we utilize the Ray distributed framework to manage concurrency.
% % For each rollout, the model will need to first generate a reasoning piece, embraced with \texttt{<think>} and \texttt{</think>}, followed by a function call embraced with \texttt{<tool\_call>} and \texttt{</tool\_call>}.
% During each rollout, the model generates a reasoning piece followed by a function call, which enables the system to handle massive parallel requests.
% When detecting the special ending token \texttt{</tool\_call>} for the function call, the model stop generation and the training framework packages the instruction alongside trajectory identifiers and image paths into batched HTTP requests. 
% Ray automatically manages load balancing and request queues on the tool server side, while our fine-grained resource control optimizes hardware utilization by binding intensive VLM tools to dedicated GPUs and multiplexing lightweight tools on CPUs.

\noindent \textbf{Extensible Tool Set.}
To support rapid expansion beyond tools used in our experiments, we provide a generic \texttt{BaseTool} interface that enables \emph{plug‑and‑play} integration of new tools with minimal boilerplate.
Operational robustness is ensured via health and metric endpoints (\texttt{/health}, \texttt{/metrics}) and Ray's automated failure recovery, which restarts crashed actors without disrupting ongoing training.

% To enable flexible and extensible tool-use capabilities for tasks that may not be covered in the experiments, the system is also designed for robustness and ease of extension.
% We introduce a generic \texttt{BaseTool} class, allowing new tools to be integrated via a standardized interface, which can rapidly expand the VLM agent's capabilities.
% Operational stability is maintained through comprehensive monitoring endpoints (\texttt{/health}, \texttt{/metrics}) and Ray’s automated failure recovery, which restarts crashed actors without disrupting the broader training process. 

